\chapter{Literature Review}
\label{ch:literature_review}

The task of automated blood vessel segmentation from medical imagery, specifically angiography and fundus photography, has been the subject of extensive academic research for several decades. The structural extraction of the vascular network is widely recognized as a critical preliminary step in the computerized analysis of a vast array of pathologies. This chapter provides a comprehensive, chronological review of the methodologies developed for blood vessel segmentation, tracing the evolution from early hand-crafted heuristic filters to the modern era of deep convolutional neural networks. Furthermore, it extensively reviews the literature surrounding the architectural innovations of the UNet and DeepLab families, the mathematical formulations designed to combat severe class imbalance, and the classical algorithms utilized for post-segmentation morphological skeletonization.

\section{Traditional and Hand-crafted Segmentation Techniques}
Prior to the proliferation of deep learning, vessel segmentation relied heavily on explicit, hand-crafted feature engineering based on the known geometric and photometric properties of blood vessels. These early methods generally assumed that vessels appear as continuous, piecewise-linear tubular structures that are darker (or lighter, depending on the imaging modality) than the surrounding tissue.

\subsection{Matched Filtering and Edge Detection}
One of the earliest and most widely adopted techniques was the matched filter response. \textcite{fraz2012blood} in their comprehensive survey noted that matched filters operate by convolving the image with multiple 2D Gaussian kernels at varying orientations. The parameters of the Gaussian curves are explicitly tuned to match the expected width profile of the blood vessels. While computationally efficient, matched filters suffer from high false-positive rates when confronted with non-vascular structures that exhibit similar contrast profiles, such as hemorrhages, exudates, or the boundaries of the optic disc.

\subsection{Hessian-based and Frangi Vesselness Filters}
To overcome the limitations of simple matched filters, researchers turned to differential geometry and multi-scale analysis. The most influential work in this domain was the formulation of the vesselness measure by Frangi et al. The Frangi filter computes the eigenvalues of the Hessian matrix at multiple spatial scales across the image. By analyzing the ratio of the eigenvalues, the filter can distinguish between blob-like structures (e.g., aneurysms), plate-like structures, and the desired tubular structures (vessels). While the Frangi filter and its numerous derivatives significantly improved the continuity of extracted vessels, they remain highly sensitive to noise and require extensive hyperparameter tuning for different datasets and imaging modalities.

\subsection{Morphological Operations}
Mathematical morphology has also played a crucial role in traditional segmentation pipelines. Operations such as the morphological top-hat transformation are frequently used to enhance vessel contrast and correct for the non-uniform background illumination inherent in clinical imaging. However, morphological methods alone are rarely sufficient for final segmentation; they are typically employed as preprocessing steps to enhance the performance of subsequent thresholding or region-growing algorithms.

The fundamental limitation of all traditional, hand-crafted methods is their inability to generalize. A filter explicitly tuned for the thick primary arteries of the optic disc will inevitably fail to detect the ultra-thin micro-capillaries in the periphery, and vice versa. The reliance on rigid heuristic rules makes these systems highly brittle in the presence of noise, artifacts, and pathological deviations.

\section{The Advent of Deep Learning in Semantic Segmentation}
The paradigm of computer vision underwent a seismic shift with the advent of deep Convolutional Neural Networks (CNNs). Unlike traditional methods, CNNs do not rely on human-engineered features; instead, they learn a hierarchy of features directly from raw pixel data through backpropagation.

\subsection{Fully Convolutional Networks (FCN)}
The breakthrough for pixel-level semantic segmentation came with the introduction of the Fully Convolutional Network (FCN) by \textcite{long2015fully}. Prior to FCN, CNNs were primarily used for image-level classification, utilizing fully connected layers that collapsed spatial dimensions into a single vector. Long et al. demonstrated that by replacing these fully connected layers with convolutional layers, the network could output a spatial map of class probabilities. Furthermore, FCN introduced the concept of transposed convolutions (or deconvolutions) to up-sample the low-resolution feature maps back to the original image dimensions. While groundbreaking, the original FCN architecture produced coarse, blob-like segmentations because the up-sampling process fundamentally struggled to recover the fine-grained spatial details lost during the successive pooling operations of the encoder.

\subsection{The U-Net Revolution}
The limitations of the FCN architecture in the medical domain were decisively addressed by \textcite{ronneberger2015u} with the introduction of the U-Net architecture. Originally developed for the segmentation of neuronal structures in electron microscopic stacks, U-Net established a symmetric, U-shaped encoder-decoder topology. The critical innovation of U-Net was the introduction of \textit{skip connections}. These connections directly concatenated the high-resolution, low-level spatial feature maps from the contracting path (encoder) with the corresponding low-resolution, high-level semantic feature maps in the expanding path (decoder). 

This structural fusion allowed the decoder to utilize both the abstract semantic context ("what is this object?") and the precise spatial localization ("where are its exact boundaries?"). U-Net rapidly became the de facto standard for medical image segmentation due to its remarkable ability to produce highly accurate segmentations even when trained on extremely small datasets, a common constraint in the medical field. The architecture inspired a massive lineage of variants, including V-Net \parencite{milletari2016v} for 3D volumetric segmentation using 3D convolutions.

\section{Advanced Architectures: UNet++ and DeepLabV3+}
Despite the overwhelming success of the standard U-Net, the segmentation of highly complex, branching, and ultra-thin structures like capillary networks revealed inherent structural limitations. The primary issue lies in the "semantic gap" between the encoder and decoder feature maps fused by the skip connections. The encoder features are structurally precise but semantically shallow, while the decoder features are semantically rich but spatially coarse.

\subsection{UNet++: Bridging the Semantic Gap}
To resolve this semantic discrepancy, \textcite{zhou2018unet++} proposed UNet++, an architecture featuring a deeply supervised, nested encoder-decoder network. Instead of a single skip connection traversing from the encoder to the decoder at each spatial level, UNet++ introduces a dense grid of intermediate convolutional blocks. 

In UNet++, the feature maps from the encoder undergo a series of dense, nested convolutions and up-sampling operations before they are finally fused with the main decoder. This intermediate processing incrementally enriches the semantic representation of the high-resolution spatial maps. By the time the encoder features reach the final decoder node, their semantic level is closely matched to the decoder features, drastically reducing the semantic gap. Furthermore, UNet++ introduced deep supervision, allowing the network to be pruned during inference by utilizing outputs from intermediate nodes, providing a trade-off between speed and accuracy. In the context of vessel segmentation, the dense skip pathways of UNet++ have been shown to significantly improve the continuity of thin vessels and reduce false-negative segmentations in low-contrast regions.

\subsection{The DeepLab Family and Atrous Convolutions}
Running parallel to the evolution of the U-Net family was the development of the DeepLab series by Chen et al. While U-Net focused on symmetric encoder-decoder structures, DeepLab focused on the mechanics of the convolution operation itself to capture multi-scale context.

Standard convolutional and pooling operations inherently reduce spatial resolution, which is detrimental for dense prediction tasks. To counteract this, the original DeepLab models introduced the concept of \textit{atrous} (or dilated) convolutions. An atrous convolution inserts "holes" (zeros) between the filter weights, effectively expanding the receptive field of the filter without increasing the number of trainable parameters and without down-sampling the feature map.

This concept was further refined in DeepLabV3 with the introduction of Atrous Spatial Pyramid Pooling (ASPP). The ASPP module probes the incoming feature map with multiple parallel atrous convolutions at different dilation rates (e.g., rates of 6, 12, and 18). This allows the network to simultaneously capture objects and context at multiple spatial scales. The outputs of these parallel branches are then concatenated and processed by a $1 \times 1$ convolution.

\textcite{chen2018encoder} subsequently introduced DeepLabV3+, which combined the multi-scale contextual power of ASPP with a simple yet effective decoder module. The encoder (often a powerful backbone like ResNet or Xception) utilizes atrous convolutions and ASPP to extract rich semantic features. The novel decoder module then up-samples these features and refines the object boundaries by fusing them with low-level features from the encoder. DeepLabV3+ achieved state-of-the-art results on benchmark semantic segmentation datasets and has seen increasing adoption in medical imaging due to its robust handling of scale variance---a critical requirement for segmenting both thick arteries and thin capillaries simultaneously.

\section{Addressing Severe Class Imbalance}
A ubiquitous challenge in the segmentation of medical structures, and particularly in angiography, is severe class imbalance. As previously noted, vessel pixels often constitute a minuscule fraction of the total image volume.

\subsection{Standard Loss and Its Limitations}
The standard loss function for binary classification is Binary Cross-Entropy (BCE). BCE evaluates the prediction of each pixel independently and averages the loss across the entire spatial map. In a heavily imbalanced scenario, BCE is dominated by the vast majority of background pixels. A network trained solely on BCE will often learn a degenerate strategy: it will simply predict "background" for every pixel, achieving a 90\%+ global accuracy while entirely failing to segment the target structure.

\subsection{Dice Loss and Soft IoU}
To combat this, the medical imaging community rapidly adopted overlap-based loss functions, most notably the Dice Loss, derived from the Sørensen-Dice coefficient. \textcite{milletari2016v} popularized the use of a differentiable "soft" Dice Loss in their V-Net paper. Unlike BCE, Dice Loss evaluates the intersection of the predicted mask and the ground truth mask relative to their union. Because it only penalizes based on the overlap of the foreground class, it is inherently invariant to the absolute number of background pixels, making it highly effective for imbalanced datasets.

\subsection{Focal Loss}
Another critical advancement in loss formulations was the introduction of Focal Loss by \textcite{lin2017focal}. Originally designed for dense object detection to address the extreme imbalance between background and foreground bounding boxes, Focal Loss is a mathematically modified version of standard Cross-Entropy. It introduces a modulating factor $(1 - p_t)^\gamma$ that dynamically scales the loss based on the confidence of the prediction. 

For easily classified pixels (e.g., clear background regions), the modulating factor approaches zero, effectively down-weighting their contribution to the total gradient. Conversely, for hard-to-classify pixels (e.g., the faint, noisy edges of micro-capillaries), the modulating factor remains high, forcing the network to focus its learning capacity on the most challenging regions of the image. Modern state-of-the-art medical segmentation pipelines frequently utilize hybrid loss functions, combining the regional overlap benefits of Dice Loss with the hard-pixel mining capabilities of BCE or Focal Loss.

\section{Post-Processing and Morphological Skeletonization}
The final output of a deep learning segmentation model is a dense probability map, which is thresholded into a binary mask. However, for many clinical applications, a binary mask is insufficient; clinicians require discrete quantitative metrics. This necessitates the use of post-processing algorithms to extract the underlying morphology of the segmented structures.

\subsection{The Skeletonization Paradigm}
Skeletonization, or morphological thinning, is a fundamental technique in digital image processing used to reduce a binary object to its structural centerline, a representation that preserves the topological and geometric properties of the original shape while discarding absolute width information. The resulting "skeleton" is ideally one pixel thick.

\subsection{The Zhang-Suen Thinning Algorithm}
One of the most classical and widely utilized thinning algorithms is the fast parallel thinning algorithm proposed by \textcite{zhang1984fast}. The Zhang-Suen algorithm operates iteratively, analyzing the $3 \times 3$ neighborhood of every foreground pixel. In each iteration, it performs two sub-iterations: one to remove south-east boundary pixels and north-west corner pixels, and another to remove north-west boundary pixels and south-east corner pixels. A pixel is deleted (turned to background) only if it satisfies specific connectivity conditions that ensure the deletion does not break the continuity of the structure or completely erode an endpoint. 

Once a vessel mask is reduced to its skeleton, subsequent algorithms can easily traverse the one-pixel-thick graph. Pixels with exactly one foreground neighbor are identified as endpoints, while pixels with three or more foreground neighbors are identified as branch points (bifurcations). By analyzing the connectivity between these nodes, algorithms can extract the exact number of distinct vessel networks, calculate the length of individual vessel segments, and map the hierarchical branching structure of the entire vascular tree.

\section{Positioning the Thesis Contribution}
While the literature on vessel segmentation is vast, a significant gap remains in the rigorous comparison of modern, high-capacity architectures specifically tailored for the unique challenges of clinical angiography. Many existing studies focus solely on retinal fundus images and evaluate simple U-Net derivatives without addressing the downstream extraction of actionable clinical metrics.

This thesis bridges that gap by directly comparing the dense skip pathways of UNet++ against the Atrous Spatial Pyramid Pooling of DeepLabV3+ on a complex angiography dataset. Furthermore, it explicitly addresses the class imbalance problem through the empirical evaluation of hybrid BCE+Dice and Focal+Dice loss formulations. Most importantly, this thesis moves beyond simple pixel-level accuracy metrics (like global accuracy or AUC) by integrating an automated mathematical skeletonization pipeline. This ensures that the output of the deep learning models is not just a visual mask, but a quantifiable set of morphological biomarkers---such as the count of distinct vessel networks---that possess direct utility for clinical diagnosis and intervention planning.
